\documentclass[11pt,a4paper]{article}
\usepackage[T1]{fontenc}
\usepackage{isabelle,isabellesym}
\usepackage{amsmath,amssymb}

% this should be the last package used
\usepackage{pdfsetup}

% urls in roman style, theory text in math-similar italics
\urlstyle{rm}
\isabellestyle{it}

%%Euro-style date: 20 September 1955
\def\today{\number\day\space\ifcase\month\or
January\or February\or March\or April\or May\or June\or
July\or August\or September\or October\or November\or December\fi
\space\number\year}

\begin{document}

\title{The Fair Games Theorem}
\author{Lawrence C. Paulson}
\maketitle

\begin{abstract}
The fair games theorem, or optional stopping theorem, gives conditions
for which the expected value of a martingale at a stopping time is equal 
to its initial expected value.
It has applications to financial mathematics 
and is \#62 on the \emph{Top 100 Matheoamtical Theorems}.
This formalisation was translated from mathlib's \textbf{OptionalStopping}
automatically by Claude~4.6. It was polished manually afterwards.
\end{abstract}

\tableofcontents

\newpage

% include generated text of all theories
\input{session}

\bibliographystyle{abbrv}

\end{document}
